%\documentclass[12pt]{report}
\documentclass[11pt, oneside]{book}
% \documentclass[11pt, oneside]{book} % Este tamaño piden en tesis
\usepackage[utf8]{inputenc} % para cargar caracteres no anglosajones 
\usepackage[spanish,english,es-tabla]{babel}
\selectlanguage{spanish}
% Alias para IEEEtran.bst cuando los campos de biblio.bib tienen language = {en}
\makeatletter
\@ifundefined{l@en}{\let\l@en\l@english}{}
\makeatother
\usepackage{datetime} % Opcional, para mejorar el formato de la fecha
\usepackage[lmargin=3cm,rmargin=2.54cm,top=2.54cm,bottom=2.54cm]{geometry}

\usepackage{etoolbox}
%%%%%% PARA UNA FUENTE QUE SE PARECE ARIAL %%%%
%\usepackage[scaled]{helvet}
%\renewcommand\familydefault{\sfdefault}  % Configura Helvetica como la fuente por defecto

\title{Redacción - Tesis - licenciatura}
\author{Juan Alonso Alcalá Lujan}
\date{\today}

% Paquetes simples y estandar 
% % % % % % % % % % % % % % % % % % % % % % % % % % %
\usepackage{bm,amsmath,amssymb,graphicx,amsfonts,amsthm}
\usepackage{fancybox} %este para variar las cajas 
\usepackage{float}
\usepackage{array}
\usepackage{comment} %%% Para comentar secciones largas
%\usepackage[justification=justified]{caption} 
\usepackage[justification=centering, singlelinecheck=false]{caption}
% Luego, reconfigura solo las figuras a justificado
%\captionsetup[figure]{justification=justified}
\usepackage{subcaption} %% para las subfiguras
\usepackage{url}
%\usepackage[numbers]{natbib} % Para usar citas narrativas y paréntetica
\usepackage{lipsum} % Para genera texto random
\usepackage[rgb]{xcolor} % Permite definir colores personalizados en formato RGB
\usepackage{booktabs} %%%%%%% Para las tablas que he exportado en python y Para tablas sin líneas verticales
\usepackage{multicol}
\usepackage{setspace}

% Compatibilidad hacia atrás con tocloft v3.0c en versiones de LaTeX kernel previas a 2024
\providecommand{\UseTaggingSocket}[2][]{}
\providecommand{\SuspendTagging}[1]{}
\providecommand{\ResumeTagging}[1]{}
\usepackage{tocloft}
\usepackage{chngcntr} % Permite cambiar el comportamiento de numeración
\usepackage{appendix}
% Redefiniciones
% % % % % % % % % % % % % % % % % % % % % % % % % % %
\renewcommand{\contentsname}{Tabla de contenidos}
\renewcommand{\bibname}{Bibliografía}
\renewcommand{\indexname}{Índice de Palabras}
\renewcommand{\figurename}{Imagen}
\renewcommand{\tablename}{TABLA}
\renewcommand{\partname}{Parte}
\renewcommand{\chaptername}{Capítulo}
\renewcommand{\appendixname}{Anexos}
\addto\captionsspanish{\renewcommand{\figurename}{FIGURA}}% Cambia "Figura" por "Imagen
\addto\captionsspanish{\renewcommand{\tablename}{TABLA}}

% Definiciones de abreviaturas de revistas astronómicas (usadas en citas BibTeX de NASA/ADS)
\providecommand{\apj}{The Astrophysical Journal}
\providecommand{\apjl}{The Astrophysical Journal Letters}
\providecommand{\apjs}{The Astrophysical Journal Supplement Series}
\providecommand{\aj}{The Astronomical Journal}
\providecommand{\aap}{Astronomy \& Astrophysics}
\providecommand{\mnras}{Monthly Notices of the Royal Astronomical Society}
\providecommand{\pasp}{Publications of the Astronomical Society of the Pacific}
\providecommand{\pasj}{Publications of the Astronomical Society of Japan}
\providecommand{\nat}{Nature}
\providecommand{\sci}{Science}
\providecommand{\araa}{Annual Review of Astronomy and Astrophysics}

% Paquetes rebuscados 
% % % % % % % % % % % % % % % % % % % % % % % % % % %
%%%%%%%%%%%% Para el abstract %%%%%%%%%%%%%%%%%%
\newenvironment{resumen}%
{\null\vfill\begin{center}%
\bfseries \fontsize{14}{16}\selectfont Resumen \end{center}}%
{\vfill\null}

\newenvironment{abstract}%
{\null\vfill\begin{center}%
\bfseries \fontsize{14}{16}\selectfont Abstract \end{center}}%
{\vfill\null}




% COMANDOS 
% % % % % % % % % % % % % % % % % % % % % % % % % % %

\providecommand{\keywords}[1] % Para palabras clave
{
  \small	
  \textbf{\textit{Palabras Clave---}} #1
}



%%%%%%%%% Cosas que no sé para que sirven %%%%%%%%%%%%%
\usepackage{units}
%\usepackage{parskip}


\newcommand{\figref}[1]{Fig. \ref{#1}} % Para modificar la forma como se escribe en la imagenes .
%\usepackage[linktocpage=true,pagebackref = true]{hyperref} 
\usepackage[linktocpage=true]{hyperref} % linktocpage Hace que los números de página en lugar del texto se vinculen en la Tabla de contenido. , %el pagebackref me servira para linkear cosas en citas de bibliografia 


\usepackage{orcidlink}
\definecolor{mycitecolor}{rgb}{0.0, 0.0, 1.0} % Ajusta estos valores para intensificar el color

\hypersetup{ 
    colorlinks=true,
    linkcolor=black, % red
    filecolor=black, % magenta      
    urlcolor=black, % red
    citecolor=black, %mycitecolor        % Color para citas de bibliografía
} %este comando de huperref me sirve para que los limks tengan diferentes colores , depende de lo que necesito

\usepackage{tikz}
\usepackage{empheq}
\usepackage{enumerate}
\usepackage{bm}
\usepackage{physics}
\usepackage{braket} %%%% Paquete que me sirve para dibujar brakets a mi antojo y facilitar cosas
% Para la forma de escribir bra-kets mirar:
% http://mirrors.ibiblio.org/CTAN/macros/latex/contrib/physics/physics.pdf

\usepackage{wasysym}
\usepackage{multirow}
\usepackage{booktabs} % Opcional, mejora la tabla

\renewcommand{\qed}{\hfill$\blacksquare$}

\usepackage{mathrsfs}
\usepackage[makeroom]{cancel}
\usepackage{setspace} % Para establecer espacios específicos

% Keywords comando creado para las palabras clave
\providecommand{\keywords}[1]
{
  \small	
  \textbf{\textit{Palabras Clave---}} #1
}


% Keywords comando creado para las palabras clave
\providecommand{\keywordsenglish}[1]
{
  \small	
  \textbf{\textit{Keywords---}} #1
}




%%%%%%%%%%%%%%%%%%%%%%%%%% SANGRÍA %%%%%%%%%%%%%%%%%%%%%%%%%%%%%%%%%%%%%%%%%%%%%%%%%%%%%%%%%
\usepackage{indentfirst}  % Asegura sangría en el primer párrafo de cada sección
\setlength{\parindent}{1.27cm}  % Define la sangría de todos los párrafos
\linespread{1.25} % Asemeja el interlineado 1.5 de Word

  
%%%%%%%%%%% Para controlar la numeración de pagina %%%%%%%%%%%%%%%%%%%%%%%%%%%%%%%%%%%%%%%%%%%%%%%%%%%%%%%%%%%%%%%%%%%%%%%%%%%%%%%%%%%%%%%%%%%%%%%%%%%%%%%%%%%%%%%%%%%%%
\usepackage{fancyhdr}
\pagestyle{fancy}
\fancyhf{}  % Limpiar encabezados y pies de página
\fancyfoot[C]{} % Borra numeración central
\fancyfoot[R]{\thepage}  % Coloca el número de página en la esquina inferior derecha

% Por defecto, LaTeX usa el estilo plain en las primeras páginas de cada capítulo. Con \fancypagestyle{plain}{...} estamos reconfigurándolo para que las páginas de inicio de capítulo sigan el mismo formato que el resto del documento.
\fancypagestyle{plain}{
    \fancyhf{} 
    \fancyfoot[R]{\thepage} % Número de página en la esquina inferior derecha
}
\renewcommand{\headrulewidth}{0pt}  % Elimina la línea en el encabezado
\renewcommand{\footrulewidth}{0pt}  % Elimina la línea en el pie de página



%%%%%%%%%%%%%%%%%%%%%%%%%%%%%%%%%%%%%%%%%%%%%%%%%%%%%%%%%%%%%%%%%%%%%%%%%%%%%%%%%%%%%%%%%%%%%%%%%%%%%%%%%%%%%%%%%%%%%%%%%%%%%%%%%%%%%%%%%%%%%%%%%%%%%%%%%%%%%%%%%%%%%%%%%%%%%%%%%%%%%%%%%%%%%%%%%%%%%%%%%%%%%%%%%%%%%%%%%%%%%%%%%%%%%%%%%%%%%%%%%%%%%%%%%%%%%%%%%%%%%%%%%%%%%%%%%%%%%%%%%%%%%%%%%%%%%%%%%%%%%%%%%%%%%%%%%%%%%%%%%%%%



%\begin{comment} % TIPOGRAFICA  
%%%%%%%%%%%% Para ver si lo puedes eliminar porque queda demasidado feo con estas caracteristicas %%%%%%%%%
% Tipo de letra: Arial.  
\usepackage{helvet}  % Usa Helvetica (similar a Arial)
\renewcommand{\familydefault}{\sfdefault}  % Configura la fuente sans-serif por defecto

% Tamaño de letra de título de capítulo: 14 puntos. 
\usepackage{titlesec}
\titleformat{\chapter}{ \centering \fontsize{14}{16}\bfseries}{\Roman{chapter}.}{1em}{}

\titleformat{\section}{\fontsize{11}{13}\bfseries}{\makebox[0.3cm][r]{\thesection}}{1em}{}

% Configurar sangría de subapartados (Tercer nivel, Números arábigos con sangría 0.8 cm)
\titleformat{\subsection}
  {\bfseries}{\makebox[0.8cm][r]{\thesubsection}~}{0pt}{}





%Tamaño de letra de texto dentro de tablas y figuras: 9 puntos.
\captionsetup[figure]{font={small,stretch=1.5}}  % small es 9pt
\captionsetup[table]{font={small,stretch=1.5}}
%\end{comment}

% Tablas y figuras: 1.50 cm. o 0.59 pulgadas. 
\captionsetup{font={stretch=1.5}}

%\setlength{\floatsep}{1.5cm}

\renewcommand{\thechapter}{\Roman{chapter}.}
\renewcommand{\thesection}{\Alph{section}.}
\renewcommand{\thesubsection}{\arabic{subsection}} % Subapartados en números arábigos (1, 2, 3...)
%\renewcommand{\thesection}{\arabic{chapter}.\arabic{section}}


\renewcommand{\footnotesize}{\fontsize{9pt}{10pt}\selectfont}  % Tamaño de fuente en notas
%\end{comment}



%\begin{comment} % INTERLINEADO
% El texto principal es a doble espacio: 2.00 cm o 0.78 pulgadas. 

\doublespacing  % Aplica doble espaciado en todo el documento

% Ecuaciones, desde 1.00 cm. hasta 3.00 cm. dependiendo de la  diagramación
\setlength{\abovedisplayskip}{1.5cm}  % Espacio antes de ecuaciones
\setlength{\belowdisplayskip}{1.5cm}  % Espacio después de ecuaciones
% \setlength{\jot}{1cm}  % Espaciado entre ecuaciones en align



% Notas de páginas: 1.00 cm. o 0.39 pulgadas. 
\setlength{\skip\footins}{1cm}  % Ajusta la distancia entre el texto y las notas al pie


%%%%%%%%%%%%%%%%%%%%%%%%%%%%%%%%%%%%%%%%%%%%%%%%%%%%%%%%%%%%%%%%%%%%%%%%%%%%%%%%%%%%%%%%%%%%%%%%%%%%%%%%%%%%%%%%%%%%%%%%%%%%%%%%%%%%%%%%%%%%%%%%%%%%%%%%%%%%%%%%%%%%%%%%%%%%%%
% numeración arabiga continua
 \counterwithout{equation}{chapter}
 \renewcommand{\theequation}{\arabic{equation}} % Asegura numeración en arábigos
\setcounter{equation}{0} % Inicia la numeración desde 1



% Evitar que las ecuaciones se reinicien en cada capítulo


% Evitar que la numeración se reinicie por capítulo
\counterwithout{table}{chapter}
\counterwithout{figure}{chapter}


%%%%%% Numeración de tablas en números romanos %%
\renewcommand{\thetable}{\Roman{table}} % Usa números romanos para tablas
% Configurar el formato de las tablas para IEEE
\captionsetup[table]{labelfont=normalfont, labelsep=newline}
%%%%%% Numeración de figuras en números romanos %%
\captionsetup[figure]{labelsep=period} % Para cambiar ":" por "."
\renewcommand{\thefigure}{\Roman{figure}}



%%%%%%%%%% Separación entre numero y texto en la listas %%%%%%%%%%%%%
\setlength{\cftfignumwidth}{7.9em}  % Ajusta la anchura del número
\setlength{\cfttabnumwidth}{6.5em}  % Ajusta la separación en la lista de tablas
% Cambiar numeración a romanos SOLO en \listoffigures y \listoftables

% Corregir la numeración en la lista de figuras (usar \cftfigaftersnum)
\renewcommand{\cftfigpresnum}{FIGURA\;} 
\renewcommand{\cftfigaftersnum}{\Roman{figure}.}

\renewcommand{\cfttabpresnum}{TABLA\;}
\renewcommand{\cfttabaftersnum}{\Roman{table}:}



% Sangría en la tabla de contenido
\setlength{\cftchapnumwidth}{0.85cm}
\setlength{\cftsecnumwidth}{0.6cm}     % Espacio reservado para el número de la sección
\setlength{\cftsecindent}{0.3cm} 
\setlength{\cftsubsecindent}{0.8cm} % Sangría de 40pt en subsecciones


\usepackage{titlesec}
\titleformat{\section}
  {\bfseries\itshape}
  {\thesection}
  {1em}
  {}
\titleformat{\subsection}
  {\bfseries\itshape}
  {\thesubsection}
  {1em}
  {}
\begin{document}


%%%%%%%% Preliminares %%%%%%%%%%%
\frontmatter
\selectlanguage{spanish}
%----------------------------------------------------------------------------------------
%	PORTADA
%----------------------------------------------------------------------------------------

\begin{titlepage}
\begin{center}
\textsc{\Large UNIVERSIDAD NACIONAL DE INGENIERÍA}\\[1em]
\textsc{\Large FACULTAD DE CIENCIAS}

% Figura del logo
\begin{figure}[ht]
\begin{center}
\includegraphics[scale=0.2]{Uni-logo_transparente_granate.png}
\end{center}
\end{figure}
\vspace{-0.5cm}
\textsc{\Large Trabajo de Investigación}\\[1em]

\begin{spacing}{1.6}
\textbf{\LARGE Reconstrucción de trayectorias de partículas cargadas
en aceleradores mediante Neural Quantum Kernels}\\[2em]
\end{spacing}

\textsc{\Large Para obtener el Grado Académico de Bachiller en Ciencias con Mención en Física}\\[1em]

% Autor con ORCID
\textsc{\large ELABORADO POR}\\[1em]
\textsc{\LARGE Hipólito Arturo Riveros Guevara} \\ \orcidlink{0009-0009-6414-2538} \textbf{\large 0009-0009-6414-2538}\\[1em] % Reemplaza con tu ORCID

% Asesor con ORCID
\textsc{\large Asesor}\\[1em]
\textsc{\LARGE Dr. Charlie Oscar Oncebay Segura} \\ \orcidlink{0000-0003-1396-8494} \textbf{\large 0000-0003-1396-8494}\\[1em]  % Reemplaza con su ORCID

\textsc{LIMA - PERÚ}\\
\textsc{2025}
\end{center}
\end{titlepage}

\setcounter{page}{2} % Portada
\begin{table}[h]
    \centering
    \renewcommand{\arraystretch}{1.3} % Espaciado entre filas
    \begin{tabular}{c l p{10cm}}
        \toprule
        \bf \centering Citar/How to cite & \multicolumn{2}{l}{Riveros Guevara [1]} \\ \cline{1-1}
        \bf \centering Referencia/Reference & [1] & \multirow{4}{10cm}{H. Riveros Guevara, \textit{``Reconstrucción de trayectorias de partículas cargadas
en aceleradores mediante Neural Quantum Kernels''} [Trabajo de Investigación]. Lima (Perú): Universidad Nacional de Ingeniería, 2025.} \\ 
        &  & \\
       \bf   Estilo/Style: & &  \\
       \bf     IEEE & &  \\
        \bottomrule
    \end{tabular}
\end{table}

\vspace{1cm} % Espacio entre tablas

\begin{table}[h]
    \centering
    \renewcommand{\arraystretch}{1.3} % Espaciado entre filas
    \begin{tabular}{cp{11cm}}
        \toprule
        \textbf{Citar/How to cite} & (Riveros, 2025)\\ \cline{1-1}
        \textbf{Referencia/Reference} &   \multirow{4}{11cm}{\hangindent=1.2cm \hangafter=1 Riveros, H. (2025). \textit{Reconstrucción de trayectorias de partículas cargadas
en aceleradores mediante Neural Quantum Kernels}. [Trabajo de Investigación, Universidad Nacional de Ingeniería]. Repositorio Institucional UNI.} \\ 
        &  \\
        \bf Estilo/Style: &  \\
        \bf APA (7ma ed.) &   \\ \bottomrule
    \end{tabular}
\end{table}

\newpage  % El como citar
%\null\vspace{\stretch{1}}
\begin{flushright}

{\bf \large Dedicatoria}

Sección opcional en la que el autor o los autores hacen mención va \\
quién va dedicada la realización del trabajo

\end{flushright}
\vspace{\stretch{2}}\null
\newpage % Dedicatoria
%\null\vspace{\stretch{1}}

\noindent {\bf \large Agradecimientos} \\

Sección opcional en la que el autor o los autores hacen mención de las personas o instituciones que contribuyeron o apoyaron en la realización del trabajo. 


\vspace{\stretch{2}}\null
\newpage % Agradecimientos
%\begin{center}
    \vspace*{3cm} % Espacio en blanco para centrar la cita en la página
    
    \begin{quote}
        \itshape
        “¿Cuánto más vas a esperar para exigir lo mejor de ti? Se te han presentado las enseñanzas esenciales y pretendes comprenderlas. Entonces, ¿qué clase de maestro estás esperando? ¿Estás retrasando la puesta en práctica de estos principios hasta que él llegue? Eres una persona adulta, no un niño. Si continúas siendo descuidado y perezoso, encontrando excusas y postergando el momento de actuar, tu falta de progreso pasará desapercibida, y al final habrás vivido y muerto sin propósito.

        En este mismo instante, pues, piensa que eres digno de vivir como un adulto que se perfecciona. Considera todo lo óptimo como una ley inviolable. Y si se te presenta un momento de dolor o de placer, de gloria o de desgracia, recuerda que el combate es ahora. Ahora es cuando comienza la Olimpíada, y no puede ser postergada.” \\
        \vspace{0.5cm}
        --- Parafraseado de Enquiridión 51, ``Epicteto"
    \end{quote}
    
    \vspace{3cm} % Espacio en blanco adicional debajo de la cita
\end{center}

\newpage




\newpage 
\phantomsection % Marca la posición exacta de la referencia
\addcontentsline{toc}{section}{Resumen}
\begin{resumen}


\noindent La física de altas energías, particularmente en el contexto del High-Luminosity Large Hadron Collider (HL-LHC), enfrenta desafíos computacionales sin precedentes en la reconstrucción de eventos. El aumento de luminosidad y del \textit{pileup} generan una explosión combinatoria en la identificación y agrupación de señales (\textit{hits}) de los detectores, poniendo a prueba las capacidades de los algoritmos clásicos. Esta limitación compromete la eficiencia y precisión de la reconstrucción, motivando la exploración de nuevos enfoques computacionales.

En este contexto, la presente tesis investiga la computación cuántica y el \textit{quantum machine learning} (QML) para la reconstrucción de trayectorias de partículas, mediante métodos de kernel cuántico, particularmente la arquitectura \textit{Neural Quantum Kernel} (NQK). Se implementa un flujo de preprocesamiento que construye \textit{tripletes} de \textit{hits} mediante cortes geométricos físicamente motivados, incorporando variantes de reducción de dimensionalidad con PCA. Sobre estos datos se desarrolla un modelo NQK basado en una red neuronal cuántica (QNN) de un qubit, con \textit{data re-uploading} y una función de costo basada en fidelidad. A partir de esta red se define un  \textit{embedding} cuántico 1-to-3, utilizado por un clasificador SVM. Se evalúan configuraciones con dimensiones de entrada \(p=9\), \(p=3\) y \(p=2\).

Los resultados muestran que la combinación de PCA y NQK mejora significativamente el rendimiento. En particular, para \(p=3\), se alcanza una precisión de 86.69 $\%$, una pureza de 88.40 $\%$ y una eficiencia de 86.26 $\%$. Estos resultados sugieren que los \textit{neural quantum kernels}, combinados con técnicas clásicas de reducción de dimensionalidad, constituyen una alternativa prometedora para una reconstrucción eficiente de trayectorias en física de altas energías.

\vspace{0.5cm}
\noindent \keywords{Aprendizaje automático cuántico, Kernels cuánticos neuronales, Reconstrucción de trayectorias de partículas, Física de altas energías.}

\end{resumen}


\newpage

\selectlanguage{english}
\phantomsection % Marca la posición exacta de la referencia
\addcontentsline{toc}{section}{Abstract}
\begin{abstract}

\noindent High-energy physics, particularly in the context of the High-Luminosity Large Hadron Collider (HL-LHC), faces unprecedented computational challenges in event reconstruction. The increase in luminosity and pileup leads to a combinatorial explosion in the identification and grouping of detector signals (hits), challenging the capabilities of current classical algorithms. This limitation compromises reconstruction efficiency and accuracy, motivating the exploration of new computational approaches.

In this context, this thesis investigates quantum computing and quantum machine learning (QML) for particle track reconstruction through quantum kernel methods, particularly the Neural Quantum Kernel (NQK) architecture. A preprocessing pipeline is implemented to construct hit triplets using physically motivated geometric cuts, incorporating dimensionality reduction variants based on Principal Component Analysis (PCA). The NQK model is built upon a single-qubit quantum neural network (QNN) with data re-uploading and a fidelity-based cost function. A 1-to-3 quantum embedding is defined from the QNN and employed by a Support Vector Machine (SVM) classifier. Configurations with input dimensions of (p=9), (p=3), and (p=2) are evaluated.

The results show that combining PCA with NQK significantly improves performance. In particular, for (p=3), the model achieves an accuracy of 86.69$\%$, a purity of 88.40$\%$, and an efficiency of 86.26$\%$. These findings suggest that neural quantum kernels, combined with classical dimensionality reduction techniques, offer a promising approach to efficient particle track reconstruction in high-energy physics.

\vspace{0.5cm}
\noindent \keywordsenglish{Quantum Machine Learning, Neural Quantum Kernels, Particle Track Reconstruction, High-Energy Physics.}
\end{abstract}

% Regresamos al español para el resto del documento
\selectlanguage{spanish}
\newpage % Resumen y Abstract 

\renewcommand{\contentsname}{\begin{center} \fontsize{14pt}{16pt}\selectfont Tabla de Contenido \end{center} \vspace{-1cm}}
\tableofcontents  % Tabla de contenidos
\newpage 


\renewcommand{\listtablename}{\begin{center} \fontsize{14pt}{16pt}\selectfont Lista de Tablas \end{center} \vspace{-1cm}}
\phantomsection
\addcontentsline{toc}{chapter}{Lista de Tablas}
\listoftables
\newpage

\renewcommand{\listfigurename}{\begin{center} \fontsize{14pt}{16pt}\selectfont Lista de Figuras \end{center} \vspace{-1cm}}
\phantomsection
\addcontentsline{toc}{chapter}{Lista de Figuras}
\listoffigures
\newpage
\renewcommand{\thesubsection}{\arabic{subsection})}


\chapter*{Introducción}
\addcontentsline{toc}{chapter}{Introducción}

%\colorbox{yellow}{\underline{\bf Introducción super general de los centros NV}}
% % % % % % % % % % % % % % % % % % % % % % % % % % % %

%\colorbox{orange}{\underline{\bf Intro centros NV} }
%%%%%%%%%%%%%%%%%%%%%%%%%%%%%%%%%%%%%%%%%%%

El estudio de sistemas binarios eclipsantes constituye una herramienta fundamental en astrofísica, ya que permite determinar parámetros físicos esenciales de las estrellas. Sin embargo, la resolución del problema inverso asociado al ajuste de modelos teóricos a datos observacionales representa un desafío computacional importante debido a la complejidad y dimensionalidad del espacio de parámetros involucrado.\\
En este trabajo se propone un modelo simplificado de sistemas binarios eclipsantes basado en el potencial de Roche, considerando componentes aproximadamente esféricos, órbitas circulares y rotación rígida, con el objetivo de reproducir adecuadamente las curvas de luz observacionales y estimar de manera confiable los parámetros físicos fundamentales del sistema. Para ello, se implementa un sistema de ajuste automático mediante algoritmos genéticos y se emplea el método de Monte Carlo para estimar las incertidumbres asociadas a los parámetros obtenidos.\\
La metodología desarrollada comprende la generación de curvas de luz sintéticas y su comparación con curvas observacionales obtenidas a partir de datos fotométricos. El grado de concordancia entre ambas se evalúa mediante el análisis de los residuos, permitiendo identificar el conjunto de parámetros que proporciona el mejor ajuste. Posteriormente, los resultados obtenidos se comparan con valores reportados en la literatura para sistemas previamente estudiados, con el fin de validar la fiabilidad del modelo y del código de optimización implementado. Finalmente, una vez validado el procedimiento, el modelo es aplicado al estudio de un sistema binario eclipsante cuyo comportamiento no ha sido caracterizado previamente, demostrando así la utilidad y el potencial del enfoque computacional propuesto para futuras investigaciones en el campo de la astrofísica estelar.


 % Introducción
%\chapter*{Objetivos}
\addcontentsline{toc}{chapter}{Objetivos}

El presente trabajo tiene como objetivo general el estudio de curvas de luz en sistemas binarios eclipsantes, abordando los siguientes aspectos fundamentales: 

\begin{enumerate}
\item Estimar eficientemente los parámetros físicos de un sistema binario eclipsante a partir de observaciones fotométricas.
\item Modelar y comparar las curvas de luz observadas y sintéticas de estrellas binarias eclipsantes.
\item Ajustar automáticamente la curva de luz sintética de un sistema binario respecto a su curva observada mediante optimización a través de un algoritmo genético.
\item Obtener las incertidumbres de los parámetros físicos del sistema por el metodo de Monte Carlo.
\item Obtener los parámetros físicos de una estrella binaria eclipsante que no han sido registrados previamente.
\end{enumerate}


%Estos objetivos buscan proporcionar una comprensión más profunda de las propiedades fotofísicas del centro NV y su aplicación en magnetometría cuántica.


%%%%%%%% Cuerpo de la tesis %%%%%%%%%%%
\mainmatter
%\chapter{Parte introductoria del trabajo}

%\input{Cuerpo_de_tesis/Capitulo4/Introducción-cap4}
%\chapter{Marco teórico y conceptual}


%
\chapter{Desarrollo del trabajo de investigación}









%\chapter{Análisis y discusión de los resultados}






%\chapter{Conclusiones}

\begin{itemize}
   \item El método empleado en este trabajo permitió estimar los parámetros físicos $i$, $q$, $\Omega_{1}$, $\Omega_{2}$ y $L$ de sistemas binarios eclipsantes a partir de datos fotométricos. Para las binarias ASAS 203113+0513.2 y ASAS 011416+0426.4, los parámetros geométricos $i$ y $q$ resultaron congruentes, dentro de sus incertidumbres, con los valores reportados en estudios previos. Sin embargo, el resto de los parámetros presentó discrepancias superiores al $10\%$ e incluso cercanas al $100\%$ en algunos casos, reflejando algunas limitaciones del método. En el caso de la binaria U Ophiuchi, la ausencia de errores en los datos fotométricos impidió calcular las incertidumbres de sus parámetros; aun así, los valores obtenidos en la etapa de optimización inicial se mantuvieron en concordancia general con los de los demás sistemas analizados.

\item Se logró modelar exitosamente las curvas de luz observacionales utilizando datos fotométricos provenientes de catálogos astronómicos, y el modelo simplificado de Roche permitió reproducir las curvas de luz sintéticas correspondientes. Los valores reducidos del indicador de ajuste SCRM evidencian una buena congruencia entre las curvas observacionales y sintéticas, confirmando que el procedimiento empleado describe correctamente el comportamiento fotométrico de los sistemas estudiados.

\item El algoritmo genético empleado cumplió su función de identificar regiones del espacio de parámetros capaces de reproducir adecuadamente las curvas de luz observacionales. Sin embargo, su carácter inherentemente aleatorio dio lugar, en algunos casos, a soluciones numéricamente diferentes en ejecuciones distintas, lo que indica que, si bien útil, no constituye el método más estable para este tipo de problemas altamente degenerados. Se concluye que los algoritmos genéticos son apropiados para una búsqueda inicial amplia, pero pueden complementarse con métodos deterministas como Simplex para alcanzar soluciones más estables y robustas.

\item El método de Monte Carlo permitió cuantificar la incertidumbre asociada a cada parámetro, incorporando la dispersión de los datos fotométricos. Las distribuciones obtenidas muestran correlaciones y regiones degeneradas del espacio de soluciones. Aunque un mayor número de iteraciones habría permitido obtener estimaciones estadísticamente más precisas y mejor caracterizadas, las limitaciones de tiempo y capacidad computacional impidieron ampliar significativamente el muestreo.

\item El sistema ASAS 061107+1619.5 fue caracterizado por primera vez en este informe. Los parámetros obtenidos presentan incertidumbres moderadas, una distribución relativamente compacta en las simulaciones de Monte Carlo y un valor reducido de SCRM  que confirma que la curva sintética reproduce satisfactoriamente la curva observacional, por lo que esta solución constituye una primera estimación fiable de la geometría del sistema. Aunque no se trata de una solución definitiva, representa un aporte significativo al conocimiento de esta fuente y sienta las bases para futuros estudios observacionales y modelados más avanzados.
\end{itemize}

 % Introducción
\chapter{Recomendaciones}

\begin{itemize}

\item Debido a la naturaleza estocástica del algoritmo genético y a la variabilidad observada entre ejecuciones, se recomienda complementar o reemplazar este método con optimizadores deterministas como Simplex, Levenberg–Marquardt o enfoques híbridos, que permiten alcanzar soluciones más estables y con menor dispersión.

\item Aunque las iteraciones realizadas permitieron obtener estimaciones razonables de las incertidumbres, un número mayor de simulaciones permitiría caracterizar con mayor precisión las distribuciones de los parámetros y reducir el impacto del ruido fotométrico.

\item El modelo simplificado basado en el potencial de Roche no contempla efectos como manchas estelares, excentricidad orbital, sincronización imperfecta o variaciones atmosféricas. Incluir estos procesos en futuros estudios permitiría reproducir con mayor fidelidad las curvas de luz observacionales y reducir discrepancias en parámetros secundarios.

\end{itemize}


%\input{secciones-de-ayuda/Preguntas-NV}

\cleardoublepage
\bibliographystyle{IEEEtran}

\chapter{Referencias bibliográficas}
\markboth{BIBLIOGRAFÍA}{BIBLIOGRAFÍA}
% Evitar título automático de \bibliography{}
\patchcmd{\thebibliography}{\chapter*{\bibname}}{}{}{} % solución muy rara pero funcional

% Incluir bibliografía sin título duplicado
\nocite{*}
\bibliography{biblio}




%%%%%%%% Anexos %%%%%%%%%%%
\backmatter

% --- Portada de Anexos (sí va al TOC) ---
\chapter*{Anexos}
\addcontentsline{toc}{chapter}{Anexos}
\vspace*{1cm}

\begin{center}
\begin{tabular}{p{0.9\textwidth}}
Anexo 1: Puntos de Lagrange \dotfill \pageref{anexo:A} \\
Anexo 2: Códigos en Python \dotfill \pageref{anexo:B} \\
\end{tabular}
\end{center}
\newpage

% --- Anexos individuales ---
\hypersetup{pageanchor=false}
\pagenumbering{arabic}
\setcounter{page}{1}

% Desactivar que figuras y tablas se añadan a las listas globales
\captionsetup{list=no}

% Configurar numeración romana continua para TODOS los anexos
%\renewcommand{\thefigure}{\Roman{figure}}
%\renewcommand{\thetable}{\Roman{table}}
%\setcounter{figure}{0}   % Reinicia UNA SOLA VEZ al empezar anexos
%\setcounter{table}{0}    % Reinicia UNA SOLA VEZ al empezar anexos

\chapter*{Anexo 1: Puntos de Lagrange}\phantomsection\label{anexo:A}
\input{Cuerpo_de_tesis/Anexos/AnexoA}

%\chapter*{Anexo 2: Demostración de la Función Flujo Relativo}\label{anexo:B}
%\input{Cuerpo_de_tesis/Anexos/AnexoB}

\chapter*{Anexo 2: Códigos en Python}\phantomsection\label{anexo:B}
\input{Cuerpo_de_tesis/Anexos/AnexoB}



\end{document}
